\section{Introduction and Falsifiable Memory Question}
\label{sec:introduction}

The problem of distinguishing fractional, delayed, and latent ecological memory effects in dynamical systems is fundamental to understanding ecosystem dynamics under perturbation.
However, as we show, this problem is well-posed only when complexity restrictions are imposed on the latent memory component.
Specifically, we assume a finite maximum latent dimension $m$, which renders the finite-horizon discrimination problem tractable.
Without such a restriction, the no-free-lunch theorem (Theorem~\ref{thm:T10}) implies that no finite-horizon experiment can distinguish between arbitrary memory hypotheses with guaranteed accuracy, rendering the problem ill-posed.

Our approach centers on a certified predator-prey model with strong Allee effect, which serves as a validated ecological nonlinear dynamics and serves as our digital twin.
We consider three competing memory hypotheses within a unified fractional-order framework:
\begin{itemize}
    \item Fractional memory (Caputo derivative of order $\alpha$),
    \item Delayed memory (discrete delay differential equation),
    \item Latent memory (finite-dimensional linear time-invariant system of dimension $m$).
\end{itemize}
All models share a common fractional order $\alpha$, two input channels (prey and predator), and are linearized around a certified coexistence equilibrium.

The core of our contribution lies in the structural separation of the fractional transfer function $s^{-\alpha}$ from both finite-dimensional rational transfer functions (latent memory) and finite-dimensional delay transfer functions.
This separation is established analytically in Theorems~\ref{thm:T4} and~\ref{thm:T5}, which show that the fractional kernel cannot be equivalently represented by any finite-dimensional integer-order latent-state model nor by any finite-delay model under the same input-output channel.

Furthermore, we leverage the finite-horizon exponential approximation of the fractional kernel (Theorem~\ref{thm:T9b}) and the associated $\mathcal{L}^1$ error bound to construct optimal excitation signals that exploit the frequency-band separation between the fractional and approximate integer-order dynamics.
The finite-horizon impossibility obstacle (Theorem~\ref{thm:T10}) guarantees that, below a certain noise floor, discrimination becomes impossible if the latent dimension is unrestricted—thus justifying our finite-dimensional assumption.

Finally, we formulate the optimal experimental design problem under energy and safety constraints, leveraging the principal-singular-vector solution for pairwise optimal inputs (Theorem~\ref{thm:T11}) and the finite-support property of the optimal frequency-domain design (Theorem~\ref{thm:T12}).
The resulting experimental protocol is validated through a reproducible simulation benchmark (Section~\ref{sec:benchmark}) and complemented by a prospective microcosm protocol for ecological validation (Section~\ref{sec:prospect}).
